\chapter{Giấc Mộng Rớt Đại Học Vẫn Thành Tỷ Phú}
\markboth{Giấc Mộng Rớt Đại Học Vẫn Thành Tỷ Phú}{Giấc Mộng Rớt Đại Học Vẫn Thành Tỷ Phú}

\begin{truyen}{Bill Gates, TikTok Và Ảo Tưởng Làm Lớn}{Startup Delusion}
\chuhoa{C}{ó bạn tuyên bố rất hùng hồn: ``Thầy ơi, Bill Gates bỏ học vẫn giàu mà. Giờ thời đại này còn có TikTok, AI, bán hàng online, crypto, đủ đường để làm lớn. Em nghĩ học đại học không phải con đường duy nhất. Có khi em nghỉ sớm để tập trung làm lớn.''}

\teacherpair{Thầy gật: ``Đúng, đại học không phải con đường duy nhất. Nhưng em đang dùng một ngoại lệ siêu hiếm để chống lại nghĩa vụ rất bình thường của mình. Bill Gates bỏ Harvard sau khi đã vào Harvard. Còn em đang đòi bỏ luôn môn Thể dục vì chạy ba vòng sân đã thở như startup hết vốn.''}{The teacher replies: ``Yes, university is not the only path. But you are using an extremely rare exception to fight a very ordinary responsibility. Bill Gates left Harvard after getting into Harvard. You are trying to quit PE because three laps already leave you breathing like a startup that ran out of funding.''}

Cả lớp cười rần, còn bạn ấy vẫn cố đỡ: ``Nhưng biết đâu em có tố chất kinh doanh?''

\teacherpair{Có thể. Nhưng tố chất không miễn cho em khỏi kỷ luật. Người ta thất bại trong khởi nghiệp không chỉ vì thiếu ý tưởng, mà vì thiếu nền tảng, thiếu bền bỉ, và quá sớm tin mình đặc biệt.}{``Maybe you do. But talent does not exempt you from discipline. Many people fail in entrepreneurship not because they lack ideas, but because they lack foundation, endurance, and because they believe too early that they are exceptional.''}

Thầy nói tiếp, giọng tỉnh rụi: ``Mạng xã hội chỉ thích kể đoạn lên nhanh, chứ nó giấu rất kỹ đoạn người ta học âm thầm nhiều năm. Em thấy một clip chốt đơn mười phút, chứ không thấy quãng thời gian người ta học sản phẩm, học giao tiếp, học chịu lỗ, học giữ chữ tín. Bỏ phần nền mà đòi lấy phần hào quang là kiểu đầu tư nhiều rủi ro nhất của tuổi trẻ.''

\textit{(The student romanticizes dropout success. The teacher distinguishes rare genius stories from ordinary self-deception.)}
\end{truyen}

\begin{baihoc}
Đừng biến câu chuyện của thiên tài thành giấy miễn học cho sự thiếu chuẩn bị của mình. Đường vòng vẫn cần năng lực thật; thời đại mới chỉ làm nó trông hào nhoáng hơn thôi.
\textit{(Do not turn a genius story into a school exemption for your own lack of preparation. Alternative paths still require real ability; the modern era only makes them look more glamorous.)}
\end{baihoc}

\begin{ghinhoanh}
\textbf{Short English to remember:}

Rare success is not a plan.

Dreams still need discipline.
\end{ghinhoanh}

\begin{tuhocnhanh}
1. Em có đang lấy người rất hiếm để tự cho phép mình bỏ phần cơ bản không? \textit{(Are you using rare people to justify skipping the basics?)}

2. Nếu thật sự có ước mơ lớn, hôm nay em đã rèn phần nền nào? \textit{(If your dream is truly big, what foundation have you trained today?)}
\end{tuhocnhanh}

\ngancach